\section{进行中的研究：RNA、可解释性与层级瓶颈}

讲座最后没有把系统包装成完成品，而是展示几项进行中的工作：直接建模原始 RNA、寻找模型内部生物特征，以及设计从患者到分子的层级信息瓶颈。它们共同说明下一代系统的难点不只是参数更多，而是怎样压缩又不丢失关键尺度。

\subsection{原始 RNA 与学习曲线}

原始 count data 噪声大、维度高、零值多，还受测序深度影响。训练曲线下降只能证明目标函数被优化，不能自动证明 learned representation 对药物发现有用。这里应把优化指标、患者外推、细胞状态恢复与真实实验命中率分开，避免把同一条 loss curve 当作所有科学价值的代理。

\lecturefigure{slide-251-raw-rna-training.jpg}{原始 RNA transcript 数据上的训练进展与曲线}{00:55:36}

图中左侧空间样本与右侧曲线把数据和优化并置。评估应同时观察 held-out likelihood、细胞类型分离、跨患者稳定性和下游实验命中率。若只优化随机 mask，模型可能擅长恢复高表达 housekeeping genes，却忽略稀有治疗通路。

\subsection{模型内部是否出现可解释生物特征}

课堂借用了前一讲“大语言模型的生物学”思路：能否在多模态癌症模型中找到对应免疫、细胞应激或上皮--间质转化的内部 feature？这不是给每个 neuron 随便命名，而是要把激活、富集、空间分布和干预联系起来。下面的单 feature 案例与 feature atlas 应被当作候选机制目录，而不是已经完成验证的生物本体。

\lecturefigure{slide-256-feature-interpretation.jpg}{从模型 feature 的高激活样本与基因集合推断生物语义}{00:56:18}

图中展示某个 feature 的 top cores、相关基因和语义摘要。证据链应包括：feature 在哪些样本激活、哪些基因贡献最大、空间上出现在哪里、在不同患者是否稳定，以及增强或抑制 feature 是否按预测改变输出。语言模型生成的标签只能作为候选解释。

\lecturefigure{slide-257-feature-atlas.jpg}{多个 learned features 对应免疫、应激、增殖与组织程序}{00:56:38}

特征图谱显示模型可能把复杂生物状态分解为可重用方向，例如 interferon response、cell proliferation 或 immune trafficking。一个 feature 可能混合多个程序，一个程序也可能分散在多个 feature 中；因此稀疏性和可读性之间存在权衡，不能把标签当作唯一真相。

\begin{warningbox}{自动标签不是机制发现}
LLM 根据 top genes 写出“干扰素反应”很方便，但它可能复述数据库常识，也可能忽略相反证据。可靠解释需要预注册的富集分析、跨 cohort 复现、空间定位和干预验证。
\end{warningbox}

\subsection{患者到分子的层级瓶颈}

最终架构必须同时处理患者、器官、肿瘤微环境、细胞邻域、单细胞和分子。若所有 token 全连接，计算不可承受；若压缩过早，稀有但决定治疗响应的信号会消失。核心问题因此从“模型要多大”转向“查询需要哪一级信息、何时汇聚、何时允许高层目标重新访问细节”。

\lecturefigure{slide-258-multimodal-bottlenecks.jpg}{大规模多模态系统的三个核心问题：瓶颈、任务与 masking policy}{00:57:54}

左侧层级从 patient 延伸到 molecules，中央问“如何构造 information bottleneck”，右侧问预测目标与 mask。工程上可采用局部编码、pooling、稀疏 attention、memory token 或检索；科学上必须检查被压缩的信息是否包含稀有细胞、克隆亚群和空间边界。

\lecturefigure{slide-259-tumor-world-model.jpg}{最终愿景：多传感器观察肿瘤，并回答患者条件治疗查询}{00:58:01}

结尾图回到起点：癌症生物学是隐藏状态，免疫荧光、H\&E 和转录组是传感器，药物是动作，问题是治疗后状态。真正闭环还需要把预测送入实验、获得新测量、更新模型，并记录失败假设。没有这个反馈循环，“world model”容易退化为静态多模态编码器。

\teachervoice{00:54:13--00:58:40，老师把 RNA 训练和 feature 分析明确称为进行中的工作，并把最终难题归结为：在哪些尺度放置 bottleneck、预测什么、遮蔽什么。讲义因此不把 demo 写成成熟临床系统。}

\begin{importantbox}{Hierarchical world model 的设计原则}
局部层学习细胞与邻域，组织层聚合空间生态，患者层整合遗传与临床条件；信息应按查询需要逐级上行，也允许高层目标向下选择细节。固定一次性 pooling 往往无法服务所有治疗问题。
\end{importantbox}

\subsection{本章小结}

进行中的研究揭示了统一系统的三个前沿：用合适分布直接建模原始 RNA，在模型内部寻找可验证的生物 feature，以及跨患者到分子尺度设计层级 bottleneck。参数规模只是资源，真正的科学问题仍是任务、信息路径和验证。

